\documentclass[
	ngerman,
	toc=listof, % Abbildungsverzeichnis sowie Tabellenverzeichnis in das Inhaltsverzeichnis aufnehmen
	toc=bibliography, % Literaturverzeichnis in das Inhaltsverzeichnis aufnehmen
	footnotes=multiple, % Trennen von direkt aufeinander folgenden Fußnoten
	parskip=half, % vertikalen Abstand zwischen Absätzen verwenden anstatt horizontale Einrückung von Folgeabsätzen
	numbers=noendperiod % Den letzten Punkt nach einer Nummerierung entfernen (nach DIN 5008)
]{scrartcl}

\usepackage[utf8]{inputenc} % muss als erstes eingebunden werden, da Meta/Packages ggfs. Sonderzeichen enthalten

% !TEX root = Projektdokumentation.tex

% Hinweis: der Titel muss zum Inhalt des Projekts passen und den zentralen Inhalt des Projekts deutlich herausstellen
\newcommand{\titel}{"Single-Sign-On"-Authentifizierungsserver mit Beispielimplementierung eines Clients}
\newcommand{\kurzertitel}{"Single-Sign-On"-Authentifizierungsserver}
\newcommand{\untertitel}{}
\newcommand{\kompletterTitel}{\titel{} -- \untertitel}

\newcommand{\autorName}{Sibylle Blümke}
\newcommand{\autorAnschrift}{Scharnweberstr. 6}
\newcommand{\autorOrt}{10247 Berlin}

\newcommand{\betriebLogo}{LogoBetrieb.png}
\newcommand{\betriebName}{wycomco GmbH}
\newcommand{\betriebAnschrift}{Fasanenstr. 35}
\newcommand{\betriebOrt}{10719 Berlin}

\newcommand{\ausbildungsberuf}{Fachinformatiker für Anwendungsentwicklung}
\newcommand{\betreff}{Dokumentation zur betrieblichen Projektarbeit}
\newcommand{\pruefungstermin}{Sommer 2018}
\newcommand{\abgabeOrt}{Berlin}
\newcommand{\abgabeTermin}{12.04.2018}
 % Metadaten zu diesem Dokument (Autor usw.)
% !TEX root = ../Projektdokumentation.tex

% Anpassung an Landessprache ---------------------------------------------------
\usepackage{babel}

% Umlaute ----------------------------------------------------------------------
%   Umlaute/Sonderzeichen wie äüöß direkt im Quelltext verwenden (CodePage).
%   Erlaubt automatische Trennung von Worten mit Umlauten.
% ------------------------------------------------------------------------------
\usepackage[T1]{fontenc}

\usepackage{textcomp} % Euro-Zeichen etc.

% Schrift ----------------------------------------------------------------------
\usepackage{lmodern} % bessere Fonts
\renewcommand*\familydefault{\sfdefault}
\usepackage[T1]{fontenc}
\usepackage{relsize} % Schriftgröße relativ festlegen

% Tabellen ---------------------------------------------------------------------
\PassOptionsToPackage{table}{xcolor}
\usepackage{tabularx}
% für lange Tabellen
\usepackage{longtable}
\usepackage{array}
\usepackage{ragged2e}
\usepackage{lscape}
\newcolumntype{w}[1]{>{\raggedleft\hspace{0pt}}p{#1}} % Spaltendefinition rechtsbündig mit definierter Breite

% Grafiken ---------------------------------------------------------------------
\usepackage{qtree}
\usepackage[dvips,final]{graphicx} % Einbinden von JPG-Grafiken ermöglichen
\usepackage{graphics} % keepaspectratio
\usepackage{floatflt} % zum Umfließen von Bildern
\graphicspath{{Bilder/}} % hier liegen die Bilder des Dokuments

% Sonstiges --------------------------------------------------------------------
\usepackage[titles]{tocloft} % Inhaltsverzeichnis DIN 5008 gerecht einrücken
\usepackage{amsmath,amsfonts} % Befehle aus AMSTeX für mathematische Symbole
\usepackage{enumitem} % anpassbare Enumerates/Itemizes
\usepackage{xspace} % sorgt dafür, dass Leerzeichen hinter parameterlosen Makros nicht als Makroendezeichen interpretiert werden

\usepackage{makeidx} % für Index-Ausgabe mit \printindex
\usepackage[printonlyused]{acronym} % es werden nur benutzte Definitionen aufgelistet

% Einfache Definition der Zeilenabstände und Seitenränder etc.
\usepackage{setspace}
\usepackage{geometry}

% Symbolverzeichnis
\usepackage[intoc]{nomencl}
\let\abbrev\nomenclature
\renewcommand{\nomname}{Abkürzungsverzeichnis}
\setlength{\nomlabelwidth}{.25\hsize}
\renewcommand{\nomlabel}[1]{#1 \dotfill}
\setlength{\nomitemsep}{-\parsep}

\usepackage{varioref} % Elegantere Verweise. „auf der nächsten Seite“
\usepackage{url} % URL verlinken, lange URLs umbrechen etc.

\usepackage{chngcntr} % fortlaufendes Durchnummerieren der Fußnoten
% \usepackage[perpage]{footmisc} % Alternative: Nummerierung der Fußnoten auf jeder Seite neu

\usepackage{ifthen} % bei der Definition eigener Befehle benötigt
\usepackage{todonotes} % definiert u.a. die Befehle \todo und \listoftodos
\usepackage[square]{natbib} % wichtig für korrekte Zitierweise

% PDF-Optionen -----------------------------------------------------------------
\usepackage{pdfpages}
\pdfminorversion=5 % erlaubt das Einfügen von pdf-Dateien bis Version 1.7, ohne eine Fehlermeldung zu werfen (keine Garantie für fehlerfreies Einbetten!)
\usepackage[
    bookmarks,
    bookmarksnumbered,
    bookmarksopen=true,
    bookmarksopenlevel=1,
    colorlinks=true,
% diese Farbdefinitionen zeichnen Links im PDF farblich aus
    linkcolor=AOBlau, % einfache interne Verknüpfungen
    anchorcolor=AOBlau,% Ankertext
    citecolor=AOBlau, % Verweise auf Literaturverzeichniseinträge im Text
    filecolor=AOBlau, % Verknüpfungen, die lokale Dateien öffnen
    menucolor=AOBlau, % Acrobat-Menüpunkte
    urlcolor=AOBlau,
% diese Farbdefinitionen sollten für den Druck verwendet werden (alles schwarz)
    %linkcolor=black, % einfache interne Verknüpfungen
    %anchorcolor=black, % Ankertext
    %citecolor=black, % Verweise auf Literaturverzeichniseinträge im Text
    %filecolor=black, % Verknüpfungen, die lokale Dateien öffnen
    %menucolor=black, % Acrobat-Menüpunkte
    %urlcolor=black,
%
    %backref, % Quellen werden zurück auf ihre Zitate verlinkt
    pdftex,
    plainpages=false, % zur korrekten Erstellung der Bookmarks
    pdfpagelabels=true, % zur korrekten Erstellung der Bookmarks
    hypertexnames=false, % zur korrekten Erstellung der Bookmarks
    linktocpage % Seitenzahlen anstatt Text im Inhaltsverzeichnis verlinken
]{hyperref}
% Befehle, die Umlaute ausgeben, führen zu Fehlern, wenn sie hyperref als Optionen übergeben werden
\hypersetup{
    pdftitle={\titel -- \untertitel},
    pdfauthor={\autorName},
    pdfcreator={\autorName},
    pdfsubject={\titel -- \untertitel},
    pdfkeywords={\titel -- \untertitel},
}


% zum Einbinden von Programmcode -----------------------------------------------
\usepackage{listings}
\usepackage{xcolor}
\definecolor{hellgelb}{rgb}{1,1,0.9}
\definecolor{colKeys}{rgb}{0,0,1}
\definecolor{colIdentifier}{rgb}{0,0,0}
\definecolor{colComments}{rgb}{0,0.5,0}
\definecolor{colString}{rgb}{1,0,0}
\lstset{
    float=hbp,
	basicstyle=\footnotesize,
    identifierstyle=\color{colIdentifier},
    keywordstyle=\color{colKeys},
    stringstyle=\color{colString},
    commentstyle=\color{colComments},
    backgroundcolor=\color{hellgelb},
    columns=flexible,
    tabsize=2,
    frame=single,
    extendedchars=true,
    showspaces=false,
    showstringspaces=false,
    numbers=left,
    numberstyle=\tiny,
    breaklines=true,
    breakautoindent=true,
	captionpos=b,
}
\lstdefinelanguage{cs}{
	sensitive=false,
	morecomment=[l]{//},
	morecomment=[s]{/*}{*/},
	morestring=[b]",
	morekeywords={
		abstract,event,new,struct,as,explicit,null,switch
		base,extern,object,this,bool,false,operator,throw,
		break,finally,out,true,byte,fixed,override,try,
		case,float,params,typeof,catch,for,private,uint,
		char,foreach,protected,ulong,checked,goto,public,unchecked,
		class,if,readonly,unsafe,const,implicit,ref,ushort,
		continue,in,return,using,decimal,int,sbyte,virtual,
		default,interface,sealed,volatile,delegate,internal,short,void,
		do,is,sizeof,while,double,lock,stackalloc,
		else,long,static,enum,namespace,string},
}
\lstdefinelanguage{natural}{
	sensitive=false,
	morecomment=[l]{/*},
	morestring=[b]",
	morestring=[b]',
	alsodigit={-,*},
	morekeywords={
		DEFINE,DATA,LOCAL,END-DEFINE,WRITE,CALLNAT,PARAMETER,USING,
		IF,NOT,END-IF,ON,*ERROR-NR,ERROR,END-ERROR,ESCAPE,ROUTINE,
		PERFORM,SUBROUTINE,END-SUBROUTINE,CONST,END-FOR,END,FOR,RESIZE,
		ARRAY,TO,BY,VALUE,RESET,COMPRESS,INTO,EQ},
}
\lstdefinelanguage{php}{
	sensitive=false,
	morecomment=[l]{/*},
	morestring=[b]",
	morestring=[b]',
	alsodigit={-,*},
	morekeywords={
		abstract,and,array,as,break,case,catch,cfunction,class,clone,const,
		continue,declare,default,do,else,elseif,enddeclare,endfor,endforeach,
		endif,endswitch,endwhile,extends,final,for,foreach,function,global,
		goto,if,implements,interface,instanceof,namespace,new,old_function,or,
		private,protected,public,static,switch,throw,try,use,var,while,xor
		die,echo,empty,exit,eval,include,include_once,isset,list,require,
		require_once,return,print,unset},
}
 % verwendete Packages
% !TEX root = ../Projektdokumentation.tex

% Seitenränder -----------------------------------------------------------------
\setlength{\topskip}{\ht\strutbox} % behebt Warnung von geometry
\geometry{a4paper,left=20mm,right=20mm,top=25mm,bottom=35mm}

\usepackage[
	automark, % Kapitelangaben in Kopfzeile automatisch erstellen
	headsepline, % Trennlinie unter Kopfzeile
	ilines % Trennlinie linksbündig ausrichten
]{scrpage2}

% Kopf- und Fußzeilen ----------------------------------------------------------
\pagestyle{scrheadings}
% chapterpagestyle gibt es nicht in scrartcl
%\renewcommand{\chapterpagestyle}{scrheadings}
\clearscrheadfoot

% Kopfzeile
\renewcommand{\headfont}{\normalfont} % Schriftform der Kopfzeile
\ihead{\large{\textsc{\kurzertitel}}\\[2ex] \textit{\headmark}}
\chead{}
\ohead{\includegraphics[scale=0.2]{\betriebLogo}}
\setlength{\headheight}{15mm} % Höhe der Kopfzeile
% \setheadwidth[0pt]{textwithmarginpar} % Kopfzeile über den Text hinaus verbreitern (falls Logo den Text überdeckt)

% Fußzeile
\ifoot{\autorName}
\cfoot{}
\ofoot{\pagemark}

% Überschriften nach DIN 5008 in einer Fluchtlinie
% ------------------------------------------------------------------------------

% Abstand zwischen Nummerierung und Überschrift definieren
% > Schön wäre hier die dynamische Berechnung des Abstandes in Abhängigkeit
% > der Verschachtelungstiefe des Inhaltsverzeichnisses
\newcommand{\headingSpace}{1.5cm}

% Abschnittsüberschriften im selben Stil wie beim Inhaltsverzeichnis einrücken
\renewcommand*{\othersectionlevelsformat}[3]{
  \makebox[\headingSpace][l]{#3\autodot}
}

% Für die Einrückung wird das Paket tocloft benötigt
%\cftsetindents{chapter}{0.0cm}{\headingSpace}
\cftsetindents{section}{0.0cm}{\headingSpace}
\cftsetindents{subsection}{0.0cm}{\headingSpace}
\cftsetindents{subsubsection}{0.0cm}{\headingSpace}
\cftsetindents{figure}{0.0cm}{\headingSpace}
\cftsetindents{table}{0.0cm}{\headingSpace}


% Allgemeines
% ------------------------------------------------------------------------------

\onehalfspacing % Zeilenabstand 1,5 Zeilen
\frenchspacing % erzeugt ein wenig mehr Platz hinter einem Punkt

% Schusterjungen und Hurenkinder vermeiden
\clubpenalty = 10000
\widowpenalty = 10000
\displaywidowpenalty = 10000

% Quellcode-Ausgabe formatieren
\lstset{numbers=left, numberstyle=\tiny, numbersep=5pt, breaklines=true}
\lstset{emph={square}, emphstyle=\color{red}, emph={[2]root,base}, emphstyle={[2]\color{blue}}}

\counterwithout{footnote}{section} % Fußnoten fortlaufend durchnummerieren
\setcounter{tocdepth}{\subsubsectionlevel} % im Inhaltsverzeichnis werden die Kapitel bis zum Level der subsubsection übernommen
\setcounter{secnumdepth}{\subsubsectionlevel} % Kapitel bis zum Level der subsubsection werden nummeriert

% Aufzählungen anpassen
\renewcommand{\labelenumi}{\arabic{enumi}.}
\renewcommand{\labelenumii}{\arabic{enumi}.\arabic{enumii}.}
\renewcommand{\labelenumiii}{\arabic{enumi}.\arabic{enumii}.\arabic{enumiii}}

% Tabellenfärbung:
\definecolor{heading}{rgb}{0.64,0.78,0.86}
\definecolor{odd}{rgb}{0.9,0.9,0.9}
 % Definitionen zum Aussehen der Seiten
% !TEX root = ../Projektdokumentation.tex

% Abkürzungen, ggfs. mit korrektem Leerraum
\newcommand{\bs}{$\backslash$\xspace}
\newcommand{\bspw}{bspw.\xspace}
\newcommand{\bzw}{bzw.\xspace}
\newcommand{\ca}{ca.\xspace}
\newcommand{\dahe}{\mbox{d.\,h.}\xspace}
\newcommand{\etc}{etc.\xspace}
\newcommand{\eur}[1]{\mbox{#1\,\texteuro}\xspace}
\newcommand{\evtl}{evtl.\xspace}
\newcommand{\ggfs}{ggfs.\xspace}
\newcommand{\Ggfs}{Ggfs.\xspace}
\newcommand{\gqq}[1]{\glqq{}#1\grqq{}}
\newcommand{\inkl}{inkl.\xspace}
\newcommand{\insb}{insb.\xspace}
\newcommand{\ua}{\mbox{u.\,a.}\xspace}
\newcommand{\usw}{usw.\xspace}
\newcommand{\Vgl}{Vgl.\xspace}
\newcommand{\zB}{\mbox{z.\,B.}\xspace}

% Befehle für häufig anfallende Aufgaben
\newcommand{\Abbildung}[1]{\autoref{fig:#1}}
\newcommand{\Anhang}[1]{\appendixname{}~\ref{#1}: \nameref{#1} \vpageref{#1}}
\newcommand{\includegraphicsKeepAspectRatio}[2]{\includegraphics[width=#2\textwidth,height=#2\textheight,keepaspectratio]{#1}}
\newcommand{\Zitat}[2][\empty]{\ifthenelse{\equal{#1}{\empty}}{\citep{#2}}{\citep[#1]{#2}}}
\newcommand{\Autor}[1]{\textsc{#1}} % zum Ausgeben von Autoren
\newcommand{\itemd}[2]{\item{\textbf{#1}}\\{#2}} % erzeugt ein Listenelement mit fetter Überschrift

% fügt Tabellen aus einer TEX-Datei ein
\newcommand{\tabelle}[3] % Parameter: caption, label, file
{\begin{table}[htbp]
\centering
\singlespacing
\input{Tabellen/#3}
\caption{#1}
\label{#2}
\end{table}}

\newcommand{\tabelleAnhang}[1] % Parameter: file
{\begin{center}
\singlespacing
\input{Tabellen/#1}
\end{center}}

% einfaches Wechseln der Schrift, z.B.: \changefont{cmss}{sbc}{n}
\newcommand{\changefont}[3]{\fontfamily{#1} \fontseries{#2} \fontshape{#3} \selectfont}

% Verwendung analog zu \includegraphics
\newlength{\myx} % Variable zum Speichern der Bildbreite
\newlength{\myy} % Variable zum Speichern der Bildhöhe
\newcommand\includegraphicstotab[2][\relax]{%
% Abspeichern der Bildabmessungen
\settowidth{\myx}{\includegraphics[{#1}]{#2}}%
\settoheight{\myy}{\includegraphics[{#1}]{#2}}%
% das eigentliche Einfügen
\parbox[c][1.1\myy][c]{\myx}{%
\includegraphics[{#1}]{#2}}%
}

\definecolor{AOBlau}{rgb}{0, 0.28, 0.56}

% verschiedene Befehle um Wörter semantisch auszuzeichnen ----------------------
\newcommand{\Index}[2][\empty]{\ifthenelse{\equal{#1}{\empty}}{\index{#2}#2}{\index{#1}#2}}
\newcommand{\Fachbegriff}[2][\empty]{\ifthenelse{\equal{#1}{\empty}}{\textit{\Index{#2}}}{\textit{\Index[#1]{#2}}}}
\newcommand{\NeuerBegriff}[2][\empty]{\ifthenelse{\equal{#1}{\empty}}{\textbf{\Index{#2}}}{\textbf{\Index[#1]{#2}}}}

\newcommand{\Ausgabe}[1]{\texttt{#1}}
\newcommand{\Eingabe}[1]{\texttt{#1}}
\newcommand{\Code}[1]{\texttt{#1}}
\newcommand{\Datei}[1]{\texttt{#1}}

\newcommand{\Assembly}[1]{\textsf{#1}}
\newcommand{\Klasse}[1]{\textsf{#1}}
\newcommand{\Methode}[1]{\textsf{#1}}
\newcommand{\Attribut}[1]{\textsf{#1}}

\newcommand{\Datentyp}[1]{\textsf{#1}}
\newcommand{\XMLElement}[1]{\textsf{#1}}
\newcommand{\Webservice}[1]{\textsf{#1}}

\newcommand{\Refactoring}[1]{\Fachbegriff{#1}}
\newcommand{\CodeSmell}[1]{\Fachbegriff{#1}}
\newcommand{\Metrik}[1]{\Fachbegriff{#1}}
\newcommand{\DesignPattern}[1]{\Fachbegriff{#1}}
 % eigene allgemeine Befehle, die z.B. die Arbeit mit LaTeX erleichtern
% !TEX root = ../Projektdokumentation.tex

% Abkürzungen, ggfs. mit korrektem Leerraum
\newcommand{\bs}{$\backslash$\xspace}
\newcommand{\bspw}{bspw.\xspace}
\newcommand{\bzw}{bzw.\xspace}
\newcommand{\ca}{ca.\xspace}
\newcommand{\dahe}{\mbox{d.\,h.}\xspace}
\newcommand{\etc}{etc.\xspace}
\newcommand{\eur}[1]{\mbox{#1\,\texteuro}\xspace}
\newcommand{\evtl}{evtl.\xspace}
\newcommand{\ggfs}{ggfs.\xspace}
\newcommand{\Ggfs}{Ggfs.\xspace}
\newcommand{\gqq}[1]{\glqq{}#1\grqq{}}
\newcommand{\inkl}{inkl.\xspace}
\newcommand{\insb}{insb.\xspace}
\newcommand{\ua}{\mbox{u.\,a.}\xspace}
\newcommand{\usw}{usw.\xspace}
\newcommand{\Vgl}{Vgl.\xspace}
\newcommand{\zB}{\mbox{z.\,B.}\xspace}

% Befehle für häufig anfallende Aufgaben
\newcommand{\Abbildung}[1]{\autoref{fig:#1}}
\newcommand{\Anhang}[1]{\appendixname{}~\ref{#1}: \nameref{#1} \vpageref{#1}}
\newcommand{\includegraphicsKeepAspectRatio}[2]{\includegraphics[width=#2\textwidth,height=#2\textheight,keepaspectratio]{#1}}
\newcommand{\Zitat}[2][\empty]{\ifthenelse{\equal{#1}{\empty}}{\citep{#2}}{\citep[#1]{#2}}}
\newcommand{\Autor}[1]{\textsc{#1}} % zum Ausgeben von Autoren
\newcommand{\itemd}[2]{\item{\textbf{#1}}\\{#2}} % erzeugt ein Listenelement mit fetter Überschrift

% fügt Tabellen aus einer TEX-Datei ein
\newcommand{\tabelle}[3] % Parameter: caption, label, file
{\begin{table}[htbp]
\centering
\singlespacing
\input{Tabellen/#3}
\caption{#1}
\label{#2}
\end{table}}

\newcommand{\tabelleAnhang}[1] % Parameter: file
{\begin{center}
\singlespacing
\input{Tabellen/#1}
\end{center}}

% einfaches Wechseln der Schrift, z.B.: \changefont{cmss}{sbc}{n}
\newcommand{\changefont}[3]{\fontfamily{#1} \fontseries{#2} \fontshape{#3} \selectfont}

% Verwendung analog zu \includegraphics
\newlength{\myx} % Variable zum Speichern der Bildbreite
\newlength{\myy} % Variable zum Speichern der Bildhöhe
\newcommand\includegraphicstotab[2][\relax]{%
% Abspeichern der Bildabmessungen
\settowidth{\myx}{\includegraphics[{#1}]{#2}}%
\settoheight{\myy}{\includegraphics[{#1}]{#2}}%
% das eigentliche Einfügen
\parbox[c][1.1\myy][c]{\myx}{%
\includegraphics[{#1}]{#2}}%
}

\definecolor{AOBlau}{rgb}{0, 0.28, 0.56}

% verschiedene Befehle um Wörter semantisch auszuzeichnen ----------------------
\newcommand{\Index}[2][\empty]{\ifthenelse{\equal{#1}{\empty}}{\index{#2}#2}{\index{#1}#2}}
\newcommand{\Fachbegriff}[2][\empty]{\ifthenelse{\equal{#1}{\empty}}{\textit{\Index{#2}}}{\textit{\Index[#1]{#2}}}}
\newcommand{\NeuerBegriff}[2][\empty]{\ifthenelse{\equal{#1}{\empty}}{\textbf{\Index{#2}}}{\textbf{\Index[#1]{#2}}}}

\newcommand{\Ausgabe}[1]{\texttt{#1}}
\newcommand{\Eingabe}[1]{\texttt{#1}}
\newcommand{\Code}[1]{\texttt{#1}}
\newcommand{\Datei}[1]{\texttt{#1}}

\newcommand{\Assembly}[1]{\textsf{#1}}
\newcommand{\Klasse}[1]{\textsf{#1}}
\newcommand{\Methode}[1]{\textsf{#1}}
\newcommand{\Attribut}[1]{\textsf{#1}}

\newcommand{\Datentyp}[1]{\textsf{#1}}
\newcommand{\XMLElement}[1]{\textsf{#1}}
\newcommand{\Webservice}[1]{\textsf{#1}}

\newcommand{\Refactoring}[1]{\Fachbegriff{#1}}
\newcommand{\CodeSmell}[1]{\Fachbegriff{#1}}
\newcommand{\Metrik}[1]{\Fachbegriff{#1}}
\newcommand{\DesignPattern}[1]{\Fachbegriff{#1}}
 % eigene projektspezifische Befehle, z.B. Abkürzungen usw.

\begin{document}
%\phantomsection
%\thispagestyle{empty}
%\pdfbookmark[1]{Eidesstattliche Erklärung}{ihkdeckblatt}
%\includegraphicsKeepAspectRatio{DeckblattIHK}{1}
%\cleardoublepage

\phantomsection
\thispagestyle{plain}
\pdfbookmark[1]{Deckblatt}{deckblatt}
% !TEX root = Projektdokumentation.tex
\begin{titlepage}

\begin{center}
\includegraphics[scale=0.25]{LogoIHK.pdf}\\[1ex]
\Large{Abschlussprüfung \pruefungstermin}\\[3ex]

\Large{\ausbildungsberuf}\\
\LARGE{\betreff}\\[4ex]

\huge{\textbf{\titel}}\\[1.5ex]
\Large{\textbf{\untertitel}}\\[8ex]

\normalsize
Abgabetermin: \abgabeOrt, den \abgabeTermin\\[3em]
\textbf{Prüfungsbewerber:}\\
\autorName\\
\autorAnschrift\\
\autorOrt\\[5ex]

\includegraphics[scale=0.4]{\betriebLogo}\\[2ex]
\textbf{Ausbildungsbetrieb:}\\
\betriebName\\
\betriebAnschrift\\
\betriebOrt\\[5em]
\end{center}

\end{titlepage}
\cleardoublepage

% Preface --------------------------------------------------------------------
\phantomsection
\pagenumbering{Roman}
\pdfbookmark[1]{Inhaltsverzeichnis}{inhalt}
\tableofcontents
\cleardoublepage

\phantomsection
\listoffigures
\cleardoublepage

\phantomsection
\listoftables
\cleardoublepage

\phantomsection
\lstlistoflistings
\cleardoublepage

\newcommand{\abkvz}{Abkürzungsverzeichnis}
\renewcommand{\nomname}{\abkvz}
\section*{\abkvz}
\markboth{\abkvz}{\abkvz}
\addcontentsline{toc}{section}{\abkvz}
% !TEX root = Projektdokumentation.tex

% Es werden nur die Abkürzungen aufgelistet, die mit \ac definiert und auch benutzt wurden. 
%
% \acro{VERSIS}{Versicherungsinformationssystem\acroextra{ (Bestandsführungssystem)}}
% Ergibt in der Liste: VERSIS Versicherungsinformationssystem (Bestandsführungssystem)
% Im Text aber: \ac{VERSIS} -> Versicherungsinformationssystem (VERSIS)

% Hinweis: allgemein bekannte Abkürzungen wie z.B. bzw. u.a. müssen nicht ins Abkürzungsverzeichnis aufgenommen werden
% Hinweis: allgemein bekannte IT-Begriffe wie Datenbank oder Programmiersprache müssen nicht erläutert werden,
%          aber ggfs. Fachbegriffe aus der Domäne des Prüflings (z.B. Versicherung)

% Die Option (in den eckigen Klammern) enthält das längste Label oder
% einen Platzhalter der die Breite der linken Spalte bestimmt.
\begin{acronym}[WWWWW]
	\acro{AD}{Active Directory}
	\acro{API}{Application Programming Interface}
	\acro{CSS}{Cascading Style Sheets}
	\acro{CSRF}{Cross-Site-Request-Forgery}
	\acro{ERM}{En\-ti\-ty-Re\-la\-tion\-ship-Mo\-dell}
	\acro{GUI}{Graphical User Interface}
	\acro{HTML}{Hypertext Markup Language}
	\acro{IDE}{Integrated Development Environment}
	\acro{JSON}{Javascript Object Notation}
	\acro{MA}{Mitarbeiter}
	\acro{MVC}{Model-View-Controller}
	\acro{MySQL}{Structured Query Language}
	\acro{ORM}{Object Relational Mapping}
	\acro{PHP}{PHP: Hypertext Preprocessor}
	\acro{REST}{Representational State Transfer}
	\acro{SAML}{Security Assertion Markup Language}
	\acro{SSO}{Single Sign On}
	\acro{SQL}{Structured Query Language}
	\acro{TDD}{Test Driven Development}
	\acro{Timy}{wycomco Zeiterfassungssoftware}
	\acro{UML}{Unified Modeling Language}
	\acro{XML}{Extensible Markup Language}
\end{acronym}

\clearpage

% Inhalt ---------------------------------------------------------------------
\pagenumbering{arabic}
% !TEX root = Projektdokumentation.tex
% !TEX root = ../Projektdokumentation.tex
\section{Einleitung}
\label{sec:Einleitung}


\subsection{Projektumfeld} 
\label{sec:Projektumfeld}

Ausbildungsbetrieb ist die \textsc{synectic software \& services gmbh}, im Folgenden als \textit{synectic} bezeichnet. \textit{synectic} ist ein mittelständisches Softwareunternehmen mit Sitz in Berlin. Zu den Produkten des Unternehmens gehören individuell anpassbare Softwarelösungen für die Arbeit im Bereich Case Management. Momentan beschäftig das Unternehmen 21 Mitarbeiter.

Der Auftraggeber des Projektes ist die Entwicklungsabteilung des Produktes \textit{syn6}. 

Ziel von \textit{syn6} ist es, Kunden neben der momentanen Desktoplösung auch eine moderne, webbasierte Lösung anzubieten.


\subsection{Projektziel} 
\label{sec:Projektziel}
Durch die Anforderungen an die Software \textit{syn6} können die Daten innerhalb der Datenbank gewissen inhaltlichen Restriktionen unterliegen. Ein triviales Beispiel hierfür wäre, dass in einem bestimmten Feld nur Daten abgespeichert werden dürfen, die mit einem Großbuchstaben beginnen. Es sind allerdings auch komplexere Regeln möglich, die mehrere Felder aus mehreren Tabellen betreffen. Beispielsweise könnte eine Bedingung lauten, dass die Summe der Textlänge der Felder $A$ und $B$ die gleiche Länge wie das Feld $C$ besitzen müssen.

Da sich das \textit{syn6} noch in einem sehr frühen Entwicklungsstadium befindet, müssen momentan viele Daten noch händisch in die Datenbank eingetragen werden. Eine Prüfroutine vor dem tatsächlichen Abspeichern der Daten in die Datenbank ist dadurch nicht möglich.

Durch das manuelle Anlegen von Daten in der Datenbank ist, gerade bei komplexeren Regeln, die Chance, dass eine dieser inhaltlichen Regeln verletzt wird relativ hoch. Solche inkonsistenten Daten erschweren die Fehlersuche, da zuerst ermittelt werden muss, ob es sich um einen Fehler in der Programmlogik handelt oder ob fehlerhafte Daten vorliegen.

Ziel ist es, ein Konsolenbasiertes Tool zu entwickeln, was anhand von vordefinierten Regeln, in der Form von \acs{SQL}-Abfragen, eine Datenbank auf inhaltliche Integrität überprüft. Diese Regeln müssen beliebig verschachtelbar sein. Das heißt, dass es möglich sein muss, die Ergebnissmengen zweier oder meherer Regeln miteinander zu verknüpfen um so eine neue Regel erstellen zu können.

Das Ergebniss einer solchen Prüfung soll in einem leicht weiterverarbeitbaren Format ausgegeben werden um es \bspw in einem Monitoring System anzeigen zu lassen.


\subsection{Projektbegründung} 
\label{sec:Projektbegruendung}
Das Hauptproblem des momentanen Validierungsprozesses ist das hohe Maß an manueller Arbeit. Da die Restriktionen teilweise sehr komplex sind und sich über mehrere Felder und Tabellen erstrecken können, ist die Chance, dass bei der manuellen Prüfung ein Fehler passiert sehr hoch. Außerdem ist dieser Prozess sehr Zeitintensiv.

Da es keinen automatisierten Weg zur Prüfung gibt, erfordern selbst kleinste Änderungen der Daten eine erneute manuelle Validierung der Datenbank. Dies kann sich massiv auf die Produktivität der Entwickler auswirken.

Des Weiteren ist es nur schwer möglich, das Ergebnis einer Prüfung weiterzuverwerten. Es existiert auch keine Möglichkeit einer Einbindung in das \ac{CI} System um die Validierung automatisch bei jedem Build auszuführen.

Diese Probleme sollen nun durch ein Tool zur automatischen Validierung der Datenbank gelöst werden. Ein weiterer Vorteil einer Softwarelösung ist es, dass sich das Tool auch problemlos auf andere Datenbanken und Projekte anwenden lässt.

\subsection{Projektschnittstellen} 
\label{sec:Projektschnittstellen}
Die Anwendung muss ausschließlich mit der zu validierenden Datenbank interagieren können. Hierfür muss der entsprechende Datenbanktreiber auf dem ausführenden System vorhanden sein.

Mit weiteren externen Systemen muss \emph{nicht} interagiert werden.

\subsection{Projektabgrenzung} 
\label{sec:Projektabgrenzung}
Das Abschlussprojekt befasst sich nicht mit dem Weiterverabeiten des Ergebnisses einer Prüfung. Die Integration in \zB ein \acs{CI}-System ist also nicht Bestandteil der Arbeit.

% !TEX root = ../Projektdokumentation.tex
\section{Projektplanung} 
\label{sec:Projektplanung}
\subsection{Projektphasen}
\label{sec:Projektphasen}
Für die Umsetzung standen 70 Stunden zur Verfügung. Diese wurden vor Projektbeginn auf die verschiedenen Phasen des Entwicklungsprozesses aufgeteilt. 
Das Ergebnis der groben Zeitplanung lässt sich dem folgendem gestapelten Balkendiagramm entnehmen. 
\begin{figure}[htb]
\centering
\includegraphicsKeepAspectRatio{timeline.png}{0.99}
\caption{grobe Zeitplanung}
\end{figure}
Eine detaillierte Übersicht der Phasen befindet sich im \Anhang{app:Zeitplanung}.
\subsection{Ressourcenplanung}
\label{sec:Ressourcenplanung}

Eine vollständige Auflistung aller während der Umsetzung des Projekts verwendeten Ressourcen befindet sich im \Anhang{app:Ressourcen}. Bei der Auswahl der verwendeten Software war es wichtig, dass hierdurch keine Zusatzkosten anfallen. Es sollte also Software verwendet werden, die entweder kostenfrei ist (\zB Open Source) oder für die wycomco bereits Lizenzen besitzt.
\subsection{Entwicklungsprozess}
\label{sec:Entwicklungsprozess}
Die Vorgehensweise nach dem das Projekt entwickelt wird, ist das Wasserfallmodell.
Bei diesem nicht iterativen, linearen Prozess werden die einzelnen Projektphasen schrittweise bearbeitet. Hierbei bilden die Phasen-Ergebnisse jeweils die bindende Vorgabe für die nächste Projektphase\footnote{vgl. \cite[S. 263]{ItHandbuch}}.
Das gesamte Projekt soll mithilfe von \ac{TDD}, zu Deutsch \textit{Testgetriebene Entwicklung}, umgesetzt werden. Das Kernprinzip von \ac{TDD} besagt, dass das Testen der Programmkomponenten den kompletten Entwicklungsprozess leitet.
Es handelt sich hierbei um eine Designstrategie in der das Testen vor der eigentlichen Implementierung stattfindet. Es soll keine Zeile Produktivcode geschrieben werden, die nicht durch einen Test vorher abgedeckt wird.
Damit lässt sich die Qualität des Codes erhöhen und den späteren Wartungsaufwand im Nachhinein zu verringern\footnote{vgl. \cite{datenschutzbeauftragter}}.
Dieses Parallelentwicklung von Code und Tests erfolgt in sich wiederholenden Mikroiterationen, die nur einen kleinen Zeitraum in Anspruch nehmen sollte. Man kann \acs{TDD} in drei Hauptteile aufspalten, die im englisch Red-Green-Refactor genannt werden. 
Die einzelnen Phasen lassen sich wie folgt beschreiben\footnote{vgl. \cite{tddwiki}}
In der ersten, der roten Phase wird ein Test geschrieben, der fehlschlägt. Dieser wird dann mit einem Minimum an Code in der grünen Phase soweit verändert, dass er erfolgreich durchläuft. In der letzten grauen Phase wird der Code refaktorisiert, dass einfacher, sauberer Quelltext entsteht. 

Zur Verdeutlichung der drei Phasen die Abbildung \ref{fig:tddcircle}\footnote{vgl. \cite{tddcircle}}
\begin{figure}[htb!]
\centering
\includegraphicsKeepAspectRatio{TDDcircle.png}{0.7}
\caption{TDD Zirkel}
\label{fig:tddcircle}%
\end{figure}
% !TEX root = ../Projektdokumentation.tex
\section{Analysephase} 
\label{sec:Analysephase}


\subsection{Ist-Analyse} 
\label{sec:IstAnalyse}

Derzeit gibt es bei der Anmeldung an den verschiedenen Diensten bei wycomco keine Optimierungen. 
Bei allen Anwendungen kann man sich nicht selbständig registrieren und sich damit ein Benutzerkonto anlegen. Der jeweilige Administrator oder Befugte muss den User per Hand per Email einladen, damit dieser die Registrierung durchführen kann und nach einem Aktivierungslink die Applikation vollständig nutzbar ist. 

Dieser Prozess lässt wie folgt beschreiben:
\begin{enumerate}
	\item Anlegen eines vorläufigen Benutzerkontos durch den Admin
	\item Benutzer bekommt eine Mail und klickt auf den Registrierungslink
	\item Benutzer füllt das Registrierungsformular aus
	\item Benutzer bekommt eine Email mit einem Aktivierungslink
	\item Benutzer klickt auf den Link und lässt sich damit aktivieren
	\item Benutzer muss sich einloggen und kann damit die Anwendung vollständig nutzen
\end{enumerate}

Wie schnell ersichtlich wird, ist der Aufwand immens, wenn dieser Vorgang auch nur dreimal wiederholt werden muss.
Dieser Prozess ist sehr zeitaufwändig. Auch im täglichen Geschäft ist es erforderlich sich an jeder Anwendung einzeln anzumelden. Wie in Abschnitt \ref{sec:Projektbegruendung} schon erwähnt, führt dies schnell zu Trivialpasswörtern und damit zu einem Sicherheitsrisiko. 


\subsection{Wirtschaftlichkeitsanalyse}
\label{sec:Wirtschaftlichkeitsanalyse}

Durch den momentanen Prozess entsteht ein hoher zeitlicher Mehraufwand, der durch die Umsetzung des Projektes verringert werden kann.
Eine Vereinheitlichung der Nutzerzugänge durch ein Single-Sign-On-Verfahren erscheint als hilfreicher Ausweg.
Damit ist es dringend nötig dieses Projekt umzusetzen. Ob sich das auch wirtschaftlich begründen lässt, wird in den nächsten Abschnitten erläutert.


\subsubsection{\gqq{Make or Buy}-Entscheidung}
\label{sec:MakeOrBuyEntscheidung}

Zu dieser Problematik gibt es eine Vielzahl von Anbietern auf dem Markt und es kommen fortlaufend Neue hinzu. 
Grundlegend für ein \ac{SSO}-Verfahren ist dessen Integrierbarkeit, ein Großteil der genutzten Anwendungen sollte unterstützt werden und dabei sollten auch Vorgaben für komplexe Passwörter und verschlüsselte Anmeldeverfahren Standard sein. Würde ein Unbefugter Zugang erhalten, hätte er in der Regel Zugriff auf alle angebundenen Anwendungen.
Auch durch ihre Nutzerfreundlichkeit sollte eine \acs{SSO} Lösung ansprechen, für Standardanwender gleichermaßen wie für Administratoren\footnote{vgl. \cite{computerwoche}}.

Ein \ac{SSO}-Verfahren, das tatsächlich alle eingesetzten Anwendungen einbinden kann und die oben genannten für wycomco erfüllt, ist kaum zu finden. Da der Großteil der entwickelten Anwendungen bei wycomco auf Laravel\footnote{vgl. \cite{Laravel}} basiert, lag es nahe die zur Verfügung gestellten Pakete zu nutzen. Dabei handelt es sich um Passport\footnote{vgl. \cite{Passport}} und Socialite\footnote{vgl. \cite{Socialite}}, die die Grundfunktionen eines SSO Servers und den angebundenen Clients bieten. Damit kann das Augenmerk auf die individuellen Anforderungen von wycomco gelegt werden und es wurde sich dazu entschieden, das Projekt in Eigenentwicklung durchzuführen.


\subsubsection{Projektkosten}
\label{sec:Projektkosten}

Im Folgenden werden die Projektkosten, die während der Entwicklung anfallen, kalkuliert. 
Dafür müssen nicht nur die Personalkosten berücksichtigt werden, sondern auch die verwendeten Ressourcen, siehe unter \ref{app:Ressourcen}. Sämtliche Werte sind Beispiel-Angaben, da im Rahmen der IHK Projektangaben auf genaue Angaben der Personalkosten verzichtet wird. 

Bei den Personalkosten wird zwischen dem Stundensatz eines Auszubildenden und eines Mitarbeiters unterschieden. Der eines Mitarbeiters wird mit \eur{40} bemessen, der eines Auszubildenden mit \eur{10}.
Für die Nutzung der Ressourcen \footnote{Hardware, Arbeitsplatz, \etc} wird ein Satz von \eur{15} angewendet.
Aus diesen Werten ergeben sich die Projektkosten in Tabelle~\ref{tab:Kostenaufstellung}.

\tabelle{Kostenaufstellung}{tab:Kostenaufstellung}{Kostenaufstellung.tex}


\subsubsection{Amortisationsdauer}
\label{sec:Amortisationsdauer}

Der Einsatz eines \ac{SSO}-Servers hat eine deutliche Zeitersparnis zur Folge. Durchschnittlich loggt sich ein Mitarbeiter pro Tag mindestens einmal am Tag in einer Anwendung ein. Bei 3 Anwendungen wie momentan bei wycomco sind es 3 Vorgänge täglich. Dazu muss er alle sechs Monat das Passwort ändern. Durch die Vielzahl an Freelancern in der Firma, schätzt man die Anzahl neuer User pro Monat auf eins. Neue User werden von den Administratoren der Anwendung eingeladen und müssen danach selbstständig ihre Registrierung vervollständigen. 

Dazu die Auflistungen in Tabelle \ref{tab:Zeitersparnis1} und Tabelle \ref{tab:Zeitersparnis2}.

\tabelle{Zeitersparnis pro Vorgang}{tab:Zeitersparnis1}{Zeitersparnis1.tex}

\tabelle{Zeitersparnis pro Monat}{tab:Zeitersparnis2}{Zeitersparnis2.tex}

Für die Zeitersparnis pro Monat ergeben sich damit 463 Minuten. 

Dies ergibt eine tägliche Ersparnis von

\begin{eqnarray}
\frac{463 \mbox{ min/Monat}}{20 \mbox{ Tage/Monat}} = 23,15 \mbox{ min/Tag}
\end{eqnarray}

Bei einer Zeiteinsparung von 23,15 Minuten pro Tag für 252 Arbeitstagen\footnote{vgl. \cite{arbeitstage}} im Jahr ergibt sich eine Zeiteinsparung von 
\begin{eqnarray}
252 \frac{Tage}{Jahr} \cdot 23,15 \frac{min}{Tag} = 5833,8 \frac{min}{Jahr} \approx 97,23 \frac{h}{Jahr} 
\end{eqnarray}

Dadurch ergibt sich eine jährliche Einsparung von 
\begin{eqnarray}
97,23 \mbox{ h} \cdot \eur{(40 + 15)}{\mbox{/h}} = \eur{5347,65}
\end{eqnarray}

Die Amortisationszeit beträgt also $\frac{\eur{2080,00}}{\eur{5347,65}\mbox{/Jahr}} \approx 0,4 \mbox{ Jahre} \approx 5 \mbox{ Monate}$.

Der Server muss also mindestens 5 Monate das alte Vorgehen ersetzen,
damit sich Anschaffungskosten und Kosteneinsparung ausgleichen. Da es vorgesehen ist die
neue Anwendung längerfristig einzusetzen, kann die Umsetzung trotz der relativ langen
Amortisationszeit auch unter wirtschaftlichen Gesichtspunkten als sinnvoll eingestuft werden.
Eine grafische Darstellung der berechneten Werte findet sich unter \ref{app:Amortisationsdiagramm} .


\subsection{Nutzwertanalyse}
\label{sec:Nutzwertanalyse}

Neben den in \fullnameref{sec:Amortisationsdauer} aufgeführten wirtschaftlichen Vorteilen ergeben sich durch Realisierung des Projekts noch einige zusätzliche Vorteile.

Wie in der \fullnameref{sec:Projektbegruendung} schon erläutert, bringt der Einsatz einer Single-Sign-On-Lösung eine höhere Sicherheit mit sich.
In einer der Anwendungen von wycomco liegen vertrauliche Kundendaten. Sollte das Projekt dort nicht umgesetzt werden und daher Unbefugte Zugriff erlangen, kann man von Opportunitätskosten sprechen. Der Fremdzugriff bringt eine Vertragsstrafe mit sich, schätzungsweise bei größeren Kunden von 60000 \eur. Das sind Kosten die vermieden werden können mit der Umsetzung von wy-connect. 

Obwohl das Projekt im Rahmen von wycomco entwickelt wurde, lässt es sich auch problemlos auf andere Kunden anwenden. Der tatsächliche Nutzen geht also über wycomco hinaus.


\subsection{Qualitätsanforderungen}
\label{sec:Qualitaetsanforderungen}

Die Qualitätsanforderungen an die Anwendung lassen sich Tabelle~\ref{tab:Qualitaetsanforderungen} entnehmen.

\tabelle{Qualitätsanforderungen}{tab:Qualitaetsanforderungen}{Qualitaetsanforderungen.tex}

\subsection{Anwendungsfälle}
\label{sec:Anwendungsfaelle}

Es wird im Zuge der Analyse des Projektes ein Anwendungsfalldiagramm erstellt. Dies stellt
Interaktionen von Benutzern mit dem System dar und zeigt somit das erwartete Verhalten der
Anwendung. Das Anwendungsfalldiagramm ist im Anhang \ref{app:Use-Case-Diagramm} und \ref{app:Use-Case-Diagramm2} dargestellt.
Bevor ein Nutzer eine Anwendung innerhalb wycomcos nutzen kann, muss ein Administrator sie einladen. Bevor der Implementierung eines \acs{SSO}-Servers geschah dies für jede Anwendung separat. Nach Einladung bekam der User eine EMail, die ihn auf das Registrierungsformular führte. Sobald er sich erfolgreich registriert hatte, musste er sich an jeder einzelnen Anwendung einloggen um sie nutzen zu können.
Mit Hilfe von wyconnect kann dieser Prozess wie in den Diagrammen gut zu sehen ist, vereinfacht werden. Administratoren erstellen einmalig eine Einladung und der Benutzer registriert und loggt sich nur einmalig ein und kann sofort alle Anwendungen nutzen.
Hierbei sei die Autorisierung des Clients vernachlässigt. 


\subsection{Lastenheft/Fachkonzept}
Am Ende der Entwurfsphase wurde zusammen mit dem Projektleiter auf Basis des Anwendungsfalldiagramms das Lastenheft erstellt. Ein Auszug befindet sich im Anhang \Anhang{app:Lastenheft}

\Zwischenstand{Analysephase}{Analyse}

% !TEX root = ../Projektdokumentation.tex
\section{Entwurfsphase} 
\label{sec:Entwurfsphase}

\subsection{Zielplattform}
\label{sec:Zielplattform}

Wie in Abschnitt \fullnameref{sec:Projektziel} erwähnt, soll am Ende des Abschlussprojektes eine eigenständige Applikation entstanden sein, die für eventuelle Weiterentwicklungen eine größtmögliche Plattform-Unabhängigkeit gewährleistet und über einen zeitgemäßen Browser oder auf mobilen Endgeräten von jedem Nutzer verwendet werden kann. 
Die Daten, auf die zugegriffen werden soll, sind in einer bestehenden MySQL Datenbank gespeichert. 
Als Programmiersprache wurde \acused{PHP}\footnote{vgl. \cite{PHP}} gewählt, da dies bereits Grundlage zahlreicher Anwendungen ist, somit leichter an wy-connect angebunden werden kann und die Entwickler wycomcos mit dieser Sprache bereits vertraut sind. 
Wie auch die bestehenden Webapplikationen wird bei der Implementation das Framework Laravel\footnote{vgl. \cite{Laravel}} genutzt.

\subsection{Framework}
\label{sec:Framework}

Laravel ist ein freies \ac{PHP}-Webframework, welches dem \ac{MVC}-Muster folgt. Es ermöglicht neben dem im folgenden unter \ref{sec:Architekturdesign} erläuterten \ac{MVC}-Architekturdesign, \acused{REST}-Webdienste zu implementieren. Laravel wird mit dem \ac{ORM} Eloquent und einem gut bedienbaren Migrationssystem ausgeliefert. Damit werden Objekte einer objektorientierten Anwendung in eine relationale Datenbank überführt.
Das Framework bringt von Hause aus ein Authentifizierungspaket mit, welches von wyconnect genutzt wird und durch die Pakete Socialite und Passport ergänzt wird, um den OAuth2 Mechanismus zu implementieren. Dazu mehr im Kapitel \fullnameref{sec:Geschaeftslogik}.

\subsection{Wahl des SSO-Protokolls}
\label{sec:Protokollwahl}

Nach der Entscheidung einen \ac{SSO}-Server einzusetzen, galt es ein geeignetes Protokoll für die Umsetzung zu finden. An erster Stelle steht dabei die Sicherheit, bei den wesentlichen Aspekten der Autorisierung und Authentifizierung. Autorisierung um bestimmten Clients den Zugriff auf die Ressourcen des Nutzers zu berechtigen und Authentifizierung um die Identität des Benutzer zu identifizieren. Eine mögliche Lösung bietet das OAuth2 Protokoll\footnote{vgl. \cite{OAuth2}}.

Die Benutzerkontrolle stellt bei OAuth2 das grundlegende Prinzip dar. Im Fokus steht der Schutz von Benutzerdaten vor unbefugtem Zugriff. Die Ressourcen von den Benutzern (als \textit{Resource Owner} bezeichnet) verwaltet ein \textit{Resource Server}, der nur autorisierte Anfragen zulässt. Ein vertrauenswürdiger \textit{Authorization Server} authentifiziert den Benutzer und holt dessen Autorisierung für den Zugriff ein. Als Autorisierungsnachweis stellt der \textit{Auhorization Server} Access Token aus. Im vorliegenden Anwendungsfall sind \textit{Resource Server} und \textit{Authorization Server} ein und derselbe Server.
Als \textit{Client} bezeichnet man dabei jede Anwendung, die auf Ressourcen eines Benutzers in seinem Namen zugreifen möchte. Im Falle der Projektarbeit handelt es sich hierbei um \textit{timy}, eine Zeiterfassungsanwendung, die die Autorin im Rahmen ihrer Ausbildung implementiert hat und die als erstes den Single-Sign-On-Dienst nutzen soll. 

Ein anderer Ansatz wäre Kerberos, deren Infrastruktur schon im \ac{AD} mit integriert ist. Allerdings ist damit die Authentifizierung auch an das \ac{AD} gebunden. Sobald sich ein Benutzer außerhalb der Domäne an eine Anwendung anmelden möchte, ist dies nicht mehr möglich. Da wycomco immer mehr Webanwendungen entwickelt und diese auch perspektivisch Kunden zur Verfügung gestellt werden sollen sollen, wäre eine Einschränkung an die AD nicht sinnvoll. 
\acs{SAML} hingegen weist eine ähnliche Struktur wie die von OAuth2 auf. Während allerdings SAML längere Nachrichten über POST Parameter senden muss, greift OAuth2 auf GET zurück. Zudem sendet SAML \acs{XML} Antworten an den Client zurück, der wesentlich komplizierter zu verarbeiten ist als das häufig genutzte \acs{JSON}, welches auch bei wycomco standardmäßig genutzt wird\footnote{vgl. \cite{SSO-Vgl}}.
OAuth2 sendet keine Benutzerinformationen mit dem AccessToken mit, wie es in SAML der Fall ist, so erhält ein Angreifer keinen direkten Zugang zu den Nutzerdaten\footnote{vgl. \cite{book}}.

Aufgrund den angeführten Merkmalen der verschiedenen Protokolle fällt die Wahl auf das OAuth2 Protokoll, welches bei wycomco zukünftig zum Einsatz kommen wird. 

\subsection{Authorization Grant Types}
\label{sec:AuthorizationGrants}

OAuth2 bringt verschiedene Arten mit um den Nutzer zu autorisieren.
\setlist[enumerate]{itemsep=0cm} 
\begin{enumerate}
\item mit dem \textbf{Passwortansatz} gibt der Nutzer direkt am Client Nutzername und Passwort ein und sendet sie per POST Request an den \ac{SSO}-Server. Das hat den Nachteil, dass der Client die Daten bearbeitet.
\item mit dem \textbf{Autorisationsansatz} lenkt der Client den Nutzer auf eine View vom \ac{SSO}-Server. Dort kann der Zugriff gestattet oder abgelehnt werden. Nach Genehmigung bekommt der Client einen Authorization Code über eine Redirect URI, die der Client beim ersten Request angegeben hat. So kann sichergestellt werden, dass der richtige Client den Code bekommt. Nun kann der Client über einen Post Request mit dem Authorization Code ein Access Token beim Server beantragen. 
\item mit dem \textbf{Implizitansatz} läuft der Prozess ähnlich dem Autorisationsansatz. Nur das senden eines Authorization Codes entfällt und der Client bekommt sofort den Access Code vom \ac{SSO}-Server geschickt.
\item mit dem \textbf{Client-Anmeldeinformationenansatz} erhält der Client ein Access Token, um auf die eigenen Ressourcen zuzugreifen und nicht im Auftrag eines Benutzers.
\end{enumerate}

Im Projekt der Autorin fallen alle Grants bis auf den zweiten heraus. Beim Passwortansatz werden die Anmeldeinformationen des Nutzers von der Anwendung bearbeitet und verschickt, dies stellt ein zu umgehendes Sicherheitsrisiko dar.
Der Client-Anmeldeinformationenansatz passt nicht zu den Anforderungen von wycomco an das Projekt, es soll eindeutig ein Nutzer authentifiziert werden, nur der Client ist nicht ausreichend. Nach Prüfung der Sicherheitsbestimmungen entschied sich die Autorin für den Autorisationsansatz. Der Nutzer kann den Zugriff gestatten und sendet seine Nutzerdaten ausschließlich über den \ac{SSO}-Server, damit erhöht sich durch die Verwendung eines Authorizationscodes die Sicherheit und der Client ist nicht in die Authentifizierung eingebunden.

\subsection{Architekturdesign}
\label{sec:Architekturdesign}

Beim Design wurde dem \ac{MVC} Muster gefolgt, welches Laravel mit sich bringt. Die Anwendung wird in 3 Komponenten aufgespaltet und sichert damit Flexibilität, Anpassbarkeit und Wiederverwendbarkeit.
Bei einer späteren Implementierung als native Anwendung, kann das Model beibehalten werden und nur die View und der Controller müssten teilweise umgeschrieben werden.

Das \textbf{Model} enthält Daten zur Weiterverarbeitung. In vielen Fällen spiegelt ein Model eine Tabelle in der Datenbank wieder, so auch beim Eloquent \ac{ORM} von Laravel. Nur in dieser Model Klasse werden die Daten bearbeitet oder erfasst. 
Die \textbf{View} ist das Benutzerinterface. Die Daten von den Models werden visualisiert und leiten \zB Benutzeraktionen und Formulare weiter. Innerhalb der View sollte ein Schreibzugriff auf Daten vermieden werden. Somit bleiben alle Daten statisch und werden nicht verändert.
Der \textbf{Controller} empfängt Anfragen (Requests) von der View. In diesen Request sind \bspw die Login Daten eines Benutzers. Der Controller verarbeitet die Daten und sendet eine Anfrage an das User Model, welches dem Controller den richtigen User zurückgibt. Dieser loggt den Benutzer ein und die View zeigt den erfolgreichen Login Prozess an. 

\subsection{Entwurf der Benutzeroberfläche}
\label{sec:Benutzeroberflaeche} 

Das Laravel Passport Paket bringt eine beispielhafte Oberfläche mit, in der alle relevanten Anforderungen Anwendung finden. Screenshots dazu sind unter \fullnameref{app:screens} zu sehen.

Über die Hauptansicht kann der Administrator die autorisierten Clients verwalten, die per \ac{SSO} auf  Benutzerdaten zugreifen dürfen. Benutzer können lediglich ihr eigenes Profil verwalten.

\subsection{Datenmodell}
\label{sec:Datenmodell}

Im folgenden sollen die wichtigsten Komponenten des Datenmodells genannt und kurz erläutert werden. Bei \fullnameref{app:ERM} werden die von der Autorin benötigten Entitäten dargestellt. 

\tabelle{Entitäten}{tab:Entitäten}{Entitaeten.tex}

\subsection{Geschäftslogik}
\label{sec:Geschaeftslogik}

Da wy-connect mit Hilfe von Laravel umgesetzt werden soll, lag es nahe bei der Implementierung auf das Passport Paket zurückzugreifen.\footnote{vgl. \cite{Passport}}
In Kapitel \fullnameref{sec:Protokollwahl} wurden bereits die verschiedenen Rollen im OAuth2 erläutert. Nachfolgend der Workflow des Protokolls.

Zuerst muss die neue App, die den Dienst nutzen möchte, registriert werden. Dabei werden Daten wie Anwendungsname und eine Weiterleitungsadresse - die \textit{redirect URI}, auf die der Nutzer weitergeleitet wird, gespeichert.

\textbf{redirect URI} -
Der Service wird den Nutzer nur zur registrierten URL weiterleiten, wodurch Angriffe durch \zB Phishing verhindert werden können. Alle redirect URIs müssen TLS verschlüsselt sein, da der Dienst nur URIs akzeptiert die mit \glqq https\grqq{}  beginnen. So soll verhindert werden, dass Token während des Autorisierungsprozesses abgefangen werden. 
\textbf{Client ID und Secret} -
Nach der Registrierung der Anwendung wird eine \textit{Client ID} und ein \textit{Client Secret}bereit gestellt. Die Client ID wird als öffentliche Information betrachtet, wird zur Erstellung von Login Seiten genutzt und stellt die öffentliche Kennung des Clients dar. 
Das Client Secret muss vertraulich behandelt werden und wird genutzt um vom Authorization Server ein Access Token zu erlangen. Das Secret ist ausschließlich der Anwendung und dem Autorisierungsserver bekannt und besteht aus einer zufälligen Zahlen-Buchstaben-Reihe. 

Der erste Schritt von OAuth 2 besteht darin, eine Autorisierung vom Benutzer zu erhalten. Dies wird dadurch erreicht, dass dem Benutzer eine vom Dienst bereitgestellte Schnittstelle angezeigt wird. Dazu wird der Nutzer über den Anmelden Button an die URL \texttt{https://wy-connect.wycomco.de/oauth/authorize} mit folgenden GET Parametern weitergeleitet:

\textbf{response\_type=code} - Gibt an, dass der Server einen Autorisierungscode erwartet \\
\textbf{client\_id} - Die Client-ID, die beim Erstellen der Anwendung festgelegt wurde \\
\textbf{redirect\_uri} - Gibt die URL an, zu dem der Benutzer nach Abschluss der Autorisierung zurückkehren soll \\
\textbf{state} - Eine zufällige Zeichenfolge, die vom Client generiert und später überprüft wird

Nachdem der Benutzer sich am Autorisierungsserver eingeloggt hat, sieht er die Autorisierungsaufforderung. Wenn der Nutzer auf \flqq Zulassen\frqq  klickt, leitet der Dienst ihn mit einem Autorisierungscode zurück zum Client. Ein Beispiel für eine Antwort wäre \texttt{https://wy-connect-client.test/login/wyconnect/callback} mit folgenden Parametern:

\textbf{code} - Der Server gibt den Autorisierungscode zurück \\
\textbf{state} - Der Server gibt den gleichen Statuswert zurück, der übergeben wurde

Dabei ist wichtig, dass der Statuswert dem vom vorherigen Request entspricht. Dieser Zustandswert wird verglichen, um sicherzustellen, dass er mit dem Ursprungswert übereinstimmt. Der Statuswert kann normalerweise in einem Cookie oder einer Sitzung gespeichert und verglichen werden, vorrausgesetzt der Benutzer kehrt zurück. Dies stellt sicher, dass der Umleitungsendpunkt nicht dazu verleitet werden kann, willkürliche Autorisierungscodes auszutauschen.

Mittels des Autorisierungscodes kann nun der Client mit folgendem Aufruf ein Access Token erlangen \texttt{POST https://wy-connect.wycomco.de/token} und den Parametern:
  
\textbf{grant\_type} = authorization\_code - Der Grant Type für diesen Ablauf ist authorization\_code \\
\textbf{code} - Dies ist der Code, den der Nutzer im vorherigen Request erhalten hat \\
\textbf{redirect\_uri} - Muss identisch mit der URL sein, die im ursprünglichen Link angegeben wurde \\
\textbf{client\_id} - Die Client-ID, die beim Erstellen der Anwendung vergeben wurde \\
\textbf{client\_secret} - Da diese Anfrage vom serverseitigen Code stammt, ist das Geheimnis enthalten

Der Server antwortet mit einem Zugriffstoken und einer Ablaufzeit. Mit diesem Access Token, kann timy dann die Nutzerdaten per API erfragen und den Nutzer einloggen\footnote{vgl. \cite{heise}}.

Im Anhang \fullnameref{app:oauth2flow} befindet sich ein Diagramm was den Ablauf des OAuth2 Authentifizierungsprozesses veranschaulicht. 

\subsection{Pflichtenheft}
\label{sec:Pflichtenheft}

Ein Beispiel für das auf dem Lastenheft (siehe Kapitel~\ref{sec:Lastenheft}: \nameref{sec:Lastenheft}) aufbauende Pflichtenheft ist im \Anhang{app:Pflichtenheft} zu finden.

\Zwischenstand{Entwurfsphase}{Entwurf}

% !TEX root = ../Projektdokumentation.tex
\section{Implementierungsphase} 
\label{sec:Implementierungsphase}

\subsection{Implementierung der Datenstrukturen}
\label{sec:ImplementierungDatenstrukturen}



\subsection{Implementierung des Kommandozeileninterfaces}
\label{sec:ImplementierungKommandozeileninterface}



\subsection{Implementierung der Geschäftslogik}
\label{sec:ImplementierungGeschaeftslogik}


\paragraph{Beispiel}


% \Zwischenstand{Implementierungsphase}{Implementierung}

% !TEX root = ../Projektdokumentation.tex
\section{Abnahme- und Einführungsphase} 
\label{sec:AbnahmeEinfuehrung}

Vor der Abnahme wurden sowohl Unittests, als auch manuelle Systemtests durchgeführt. Ein Screenshot der Testsuite findet sich im Anhang. (Anhang noch einfügen)
Alle Tests  lieferten ein positives Ergebnis und das Projekt konnte über den firmenintern Deployment Prozess ausgerollt werden.
Abschließend wurde der Code von der Clientanwendung timy in den master Branch des Projektes gemerged und damit automatisiert auf den Webserver gespielt.
Damit konnte produktiv der SSO-Server mit der ersten Client Anwendung laufen. 
Diese neue Funktionalität wurde vor dem Teamleiter präsentiert, der anschließend das Projekt abnahm.
\Zwischenstand{Abnahme- und Einführungsphase}{Einfuehrung}

% !TEX root = ../Projektdokumentation.tex
\section{Dokumentation}
\label{sec:Dokumentation}

\begin{itemize}
	\item Wie wurde die Anwendung für die Benutzer/Administratoren/Entwickler dokumentiert (\zB Benutzerhandbuch, \acs{API}-Dokumentation)?
	\item Hinweis: Je nach Zielgruppe gelten bestimmte Anforderungen für die Dokumentation (\zB keine IT-Fachbegriffe in einer Anwenderdokumentation verwenden, aber auf jeden Fall in einer Dokumentation für den IT-Bereich).
\end{itemize}

\paragraph{Beispiel}
Ein Ausschnitt aus der erstellten Benutzerdokumentation befindet sich im 
%\Anhang{app:BenutzerDoku}.
....
Die Entwicklerdokumentation wurde mittels PHPDoc\footnote{Vgl. \cite{phpDoc}} automatisch generiert. Ein beispielhafter Auszug aus der Dokumentation einer Klasse findet sich im
% \Anhang{app:Doc}. 

\Zwischenstand{Dokumentation}{Dokumentation}

% !TEX root = ../Projektdokumentation.tex
\section{Fazit} 
\label{sec:Fazit}

\subsection{Soll-/Ist-Vergleich}
\label{sec:SollIstVergleich}
Rückblickend betrachtet, kann festgehalten werden, dass alle Anforderungen gemäß dem Pflichtenheft erfüllt wurden. 
In Tabelle~\ref{tab:Vergleich} wird die Zeit, die tatsächlich für die einzelnen Phasen benötigt wurde, der zuvor eingeplanten Zeit gegenübergestellt. 
Es ist zu erkennen, dass nur geringfügig von der Zeitplanung abgewichen wurde. Die sich daraus ergebenen Differenzen konnten untereinander kompensiert werden, sodass das Projekt in dem von der IHK festgelegten Zeitrahmen von 70 Stunden umgesetzt werden konnte.
Die zeitlichen Abweichungen kamen auf Grund von Mehraufwand bei der Recherche des OAuth2 Protokolls und gemindertem Aufwand bei der Implementierung auf Grund des verwendeten Frameworks zu Stande.
\tabelle{Soll-/Ist-Vergleich}{tab:Vergleich}{Zeitnachher.tex}

\subsection{Lessons Learned}
\label{sec:LessonsLearned}
Im Zuge der Projektdurchführung konnte die Autorin umfangreiche Erfahrungen auf dem Gebiet der Planung und Umsetzung sammeln.
Besonders deutlich wurde wie wichtig eine gute Recherche ist, da ohne gründliches Wissen über das verwendete OAuth2 Protokoll keine Implementierung möglich gewesen wäre.
Auch das verwendete Framework Laravel und seine Pakete erleichterten zwar die eigentliche Implementierung, aber ohne Verständnis und vorherige Recherche kann dieses Werkzeug nicht bedient werden. 
\subsection{Ausblick}
\label{sec:Ausblick}
Das Projekt findet aktuell Anwendung in den ersten Clientanwendungen und soll künftig noch auf Dienste ausgeweitet werden, die nicht mit Laravel geschrieben wurden und somit noch angepasst werden müssen. Auch ist es denkbar die Sicherheit zu erhöhen und eine Zwei-Faktor-Authentifizierung, oder eine Anbindung an die AD einzuführen.


% Literatur ------------------------------------------------------------------
\clearpage
\renewcommand{\refname}{Literaturverzeichnis}
\bibliography{Bibliographie}
\bibliographystyle{Allgemein/natdin} % DIN-Stil des Literaturverzeichnisses
% !TEX root = Projektdokumentation.tex
\clearpage
\addsec{Eidesstattliche Erklärung}

% Hinweis: die eidesstattliche Erklärung ist ggfs. an die Richtlinie der IHK anzupassen

Ich, \autorName, versichere hiermit eidesstattlich, dass ich meine \textbf{\betreff} mit dem
Thema
\begin{quote}
\textit{\kompletterTitel}
\end{quote}
selbständig verfasst und keine anderen als die angegebenen Quellen und Hilfsmittel benutzt habe,
wobei ich alle wörtlichen und sinngemäßen Zitate als solche gekennzeichnet habe. Die Arbeit
wurde bisher keiner anderen Prüfungsbehörde vorgelegt und auch nicht veröffentlicht.\\[6ex]

\abgabeOrt, den \abgabeTermin


\rule[-0.2cm]{5.5cm}{0.5pt}

\textsc{\autorName}


% Anhang ---------------------------------------------------------------------
\clearpage
\appendix
\pagenumbering{roman}
% !TEX root = Projektdokumentation.tex
\section{Anhang}
%\subsection{Detaillierte Zeitplanung}
%\label{app:Zeitplanung}

%\tabelleAnhang{ZeitplanungKomplett}

\subsection{Verwendete Resourcen}
\label{app:Ressourcen}

\subsubsection{Hardware}

\begin{itemize}
	\item Büroarbeitsplatz mit privatem MacBook Pro
\end{itemize}

\subsubsection{Software}

\begin{itemize}
	\item OS X 10.11 El Capitan – Betriebssystem
	\item StarUML – Anwendung zum Erstellen von \acs{UML}-Diagrammen
	\item Scala – Programmiersprache
	\item IntelliJ Community Edition – Entwicklungsumgebung Scala
	\item git – Verteilte Versionsverwaltung
	\item Gitlab – Selfhosted Repository Verwaltung
	\item texmaker – \LaTeX\xspace Editor
	\item neovim – Moderne Distribution des Editors \textit{vim}
	\item Postgresql – Datenbanksystem
	\item pgAdmin – Verwaltungswerkzeug für Postgres Datenbanken
\end{itemize}

\newpage
\subsection{Lastenheft (Auszug)}
\label{app:Lastenheft}
Es folgt ein Auszug aus dem Lastenheft mit Fokus auf die Anforderungen:

Die Anwendung muss folgende Anforderungen erfüllen: 
\begin{enumerate}[itemsep=0em,partopsep=0em,parsep=0em,topsep=0em]
\item Single-Sign-On-Server
	\begin{enumerate}
	\item Die Anwendung verfügt über eine eigene Datenbank für die Benutzerdaten.
	\item Anzeigen einer Übersichtsseite für autorisierte Clienten und Token mit allen relevanten Informationen zu diesen.
	\item Die Authentifizierung muss über einen sicheren Kanal erfolgen.
	\end{enumerate}
\item Login Prozess der Client Anwendung
	\begin{enumerate}
	\item Die Client Anwendung verfügt über eine eigene Benutzerdatenbank.
	\item Die Authentifizierungsmöglichkeit über wy-connect muss einfach auf andere Clients übertragbar sein.
	\item Der Nutzer muss die Client Anwendung zum Zugriff autorisieren.
	\item Der Nutzer meldet sich nur am Server mit seinen Nutzerdaten an.
	\end{enumerate}
\item Sonstige Anforderungen
	\begin{enumerate}
	\item Der Server sowie der Login Prozess am Client soll über eine \acs{GUI} intuitiv bedient werden können.
	\item Der Server sowie der Client müssen ohne Installation von zusätzlicher Software über einen
Webbrowser im Intranet erreichbar sein. 
	\item Zur Versionskontrolle der Anwendungsentwicklung soll ein Git\footnote{vgl. \cite{Git}}-Repository
verwendet werden.
	\item Die Anwendung soll in der Programmiersprache \acs{PHP}\footnote{vgl. \cite{PHP}} mittels des Frameworks Laravel\footnote{vgl. \cite{Laravel}} umgesetzt werden.	
	\item Der Server soll auf dem firmeninternen Webserver gehostet
werden.
	\item Bei Einsatz einer MySQL-Datenbank\footnote{vgl. \cite{MySQL}} soll der firmeninterne MySQL Server Verwendung
finden.
	\item Der Einrichtungsprozess für neue Clienten muss dokumentiert werden.
	\end{enumerate}
\end{enumerate}
		{[...]}



\clearpage

\subsection{Use-Case-Diagramm ohne SSO}
\label{app:Use-Case-Diagramm}
Das folgende Diagramm beschreibt den Anmeldeprozess an verschiedenen Systemen ohne \acs{SSO}.
\begin{figure}[htb]
\centering
\includegraphicsKeepAspectRatio{use_case_diagramm.png}{0.9}
\caption{Use-Case-Diagramm ohne SSO}
\end{figure}
\clearpage

\subsection{Use-Case-Diagramm mit SSO}
\label{app:Use-Case-Diagramm2}
Das folgende Diagramm beschreibt den Anmeldeprozess an verschiedenen Systemen mit \acs{SSO}.
\begin{figure}[htb]
\centering
\includegraphicsKeepAspectRatio{use_case_diagramm2.png}{0.9}
\caption{Use-Case-Diagramm mit SSO}
\end{figure}
\clearpage

\subsection{Benutzeroberfläche von wyconnect und dem Client}
\label{app:wyconnect-mockup}

Das folgende Template beschreibt die Benutzeroberfläche an verschiedenen Systemen mit \acs{SSO}.
\begin{figure}[htb]
\centering
\includegraphicsKeepAspectRatio{mockup-server.png}{0.9}
\caption{Benutzeroberfläche des SSO-Servers}
\end{figure}
\clearpage

Das folgende Template beschreibt den Anmeldeprozess am \acs{SSO}-Server.
\begin{figure}[htb]
\centering
\includegraphicsKeepAspectRatio{login-wyconnect.png}{0.9}
\caption{Login am SSO-Servers}
\end{figure}
\clearpage

Das folgende Mockup beschreibt den Anmeldeprozess am Client mit \acs{SSO}.
\begin{figure}[htb]
\centering
\includegraphicsKeepAspectRatio{mockup-timy.png}{0.9}
\caption{Login am Client Timy}
\end{figure}
\clearpage

\subsection{Entity Relationship Modell}
\label{app:ERM}
Das folgende Diagramm zeigt die Entitäten und Attribute der Datenbank und deren Beziehung zueinander.
\begin{figure}[htb]
\centering
\includegraphicsKeepAspectRatio{ERM.png}{0.99}
\caption{ERM}
\end{figure}
\clearpage

\subsection{Amortisationsdiagramm}
\label{app:Amortisationsdiagramm}
Das folgende Diagramm zeigt die Amortisation
\begin{figure}[htb]
\centering
\includegraphicsKeepAspectRatio{amortisation.png}{0.9}
\caption{Amortisations-Diagramm}
\end{figure}
\clearpage

\subsection{Sequenzdiagramm OAuth2}
\label{app:sequenzdiagramm}
Das folgende Diagramm zeigt den Ablauf des OAuth2 Protokolls am Beispiel des Clients Timy.
\begin{figure}[htb]
\centering
\includegraphicsKeepAspectRatio{sequenzdiagrammSSO.jpg}{0.99}
\caption{Sequenzdiagramm OAuth2}
\end{figure}
\clearpage

\subsection{Sequenzdiagramm OAuth2}
\label{app:sequenzdiagramm}
Das folgende Diagramm zeigt den Ablauf des OAuth2 Protokolls am Beispiel des Clients Timy.
\begin{figure}[htb]
\centering
\includegraphicsKeepAspectRatio{sequenzdiagrammSSO.jpg}{0.99}
\caption{Sequenzdiagramm OAuth2}
\end{figure}
\clearpage

\newpage
\subsection{Pflichtenheft (Auszug)}
\label{app:Pflichtenheft}

Die geplante Umsetzung der im Lastenheft (Auszug siehe \ref{app:Lastenheft}) definierten Anforderungen wird
in folgendem Auszug aus dem Pflichtenheft beschrieben:

\subsubsection*{Umsetzung der Anforderungen}

\begin{enumerate}[itemsep=0em,partopsep=0em,parsep=0em,topsep=0em]

\item Single-Sign-On-Server
	\begin{enumerate}
		\item Für das speichern der Benutzerdaten wird die firmeninterne MySQL Datenbank genutzt. 
		\item Innerhalb des \acs{SSO-Servers} gibt eine View für die Benutzerverwaltung:	
		\begin{itemize}
			\item Admins können Nutzer anlegen, löschen, bearbeiten
			\item Benutzer können nur ihr eigenes Profil löschen und bearbeiten
		\end{itemize}
	\item Als Authentifizierungsprotokoll wird OAuth2 \footnote{vgl. \cite{OAuth2}}genutzt.
	\begin{itemize}
		\item die Autorisierung wird hierbei über die API vorgenommen
		\item die Authentifizierung erfolgt ausschließlich über den Server
		\item zur Umsetzung wird das Laravel Paket Passport\footnote{vgl. \cite{Passport}} genutzt
	\end{itemize}
	\end{enumerate}

\item Login Prozess der Client Anwendung
		\begin{enumerate}
		\item Für das speichern der Benutzerdaten wird die firmeninterne MySQL Datenbank genutzt. 
		\item Die wy-connect Anbindung an andere Laravel Anwendungen wird mittels Installation eines bereitgestellten Pakets umgesetzt.
		\begin{itemize}
		\item dieses Paket kann über einen einfachen Konsolenbefehl in den Clienten integriert werden
		\item der wy-connect Provider liegt dabei auf der firmeneigenen Gitlab Instanz
		\item zur Umsetzung wird das Laravel Paket Socialite\footnote{vgl. \cite{Socialite}} genutzt
	\end{itemize}
		\item Nachdem Login Prozess am \acs{SSO} wird der Benutzer gefragt ob der Test Client auf die Nutzerdaten zugreifen darf.
		\item Innerhalb des Login Prozesses am Clienten werden die Nutzerdatzen nur vom \acs{SSO}-Server abgefragt.
		\end{enumerate}
\item Sonstige Anforderungen
		\begin{enumerate}
		\item Über das Webinterface kann sich ein Nutzer anmelden und hat über eine übersichtlich gestaltete  \acs{GUI} Zugriff auf seine Einstellungen.
		\item Das Programm läuft als Webanwendung.
		\item Zur Versionskontrolle wird der firmeneigene GitLab-Server\footnote{vgl. \cite{GitLab}} genutzt.\\
		{[...]}
		\end{enumerate}
\end{enumerate}




\end{document}
